\chapter{Uric Acid}

\section*{What Is This Test?}
The uric acid blood test measures the concentration of uric acid (urate) in the serum. Uric acid is the final breakdown product of purine metabolism in humans. Purines (adenine and guanine) from DNA/RNA turnover and dietary sources are metabolized through a series of enzymatic steps: purine $\rightarrow$ hypoxanthine $\rightarrow$ xanthine $\rightarrow$ uric acid, with the final step catalyzed by xanthine oxidase. Unlike most mammals, humans lack the enzyme uricase (urate oxidase), which converts uric acid to the more soluble allantoin. This evolutionary loss makes humans uniquely susceptible to hyperuricemia and gout. At physiological pH 7.4, approximately 98\% of uric acid exists as monosodium urate, which has limited solubility in plasma (approximately 6.8 mg/dL at 37 degrees C). When the concentration exceeds this solubility limit, monosodium urate crystals can precipitate in joints, soft tissues, and the kidneys. Uric acid is excreted approximately two-thirds by the kidneys (filtration, near-complete reabsorption in the proximal tubule, followed by secretion and partial post-secretory reabsorption via URAT1 transporter) and one-third by the GI tract (degraded by gut bacteria). Serum uric acid level reflects the balance between purine production and uric acid excretion. Hyperuricemia is defined as serum urate >6.8 mg/dL, the solubility limit, though gout typically requires sustained levels above this threshold for years.

\section*{Alternative Names}
Urate; Serum Uric Acid; Uric Acid, Serum; SUA; Blood Urate; Uricemia

\section*{Principle}
Serum uric acid is measured by the uricase enzymatic method: uric acid + O\textsubscript{2} + 2 H\textsubscript{2}O $\rightarrow$ allantoin + CO\textsubscript{2} + H\textsubscript{2}O\textsubscript{2} (uricase). The reaction product is detected by one of two approaches: (1) UV spectrophotometric method: uric acid absorbs light at 293 nm, and the decrease in absorbance as uric acid is consumed is proportional to uric acid concentration; or (2) coupled peroxidase method: H\textsubscript{2}O\textsubscript{2} + chromogen $\rightarrow$ colored product (measured at 500--540 nm). The coupled peroxidase method is the most common in automated chemistry analyzers. Colorimetric methods using phosphotungstic acid (Folin-Wu method) were historically used but are non-specific and have been replaced by the enzymatic method. HPLC is the reference method. The test is performed on serum separator tubes. The spectrophotometric uricase method is highly specific and precise, with minimal interference from other serum constituents. This test is NOT performed by hematology analyzers.

\section*{History}
Uric acid was first isolated from kidney stones (hence the name ``uric'' from Greek ``ouron'' for urine) by Carl Wilhelm Scheele in 1776. The relationship between uric acid and gout was established by Alfred Baring Garrod in 1848, who demonstrated uric acid in the blood of gout patients (the first recognition of a biochemical abnormality in disease). Xanthine oxidase was discovered in 1902. Allopurinol, the first xanthine oxidase inhibitor, was developed by Gertrude Elion and George Hitchings in the 1960s, for which they received the Nobel Prize in 1988. Febuxostat, a non-purine xanthine oxidase inhibitor, was approved in 2009. The URAT1 transporter and its role in urate handling was discovered in 2002, leading to the development of uricosuric agents like lesinurad.

\section*{Why This Test Is Ordered}
\begin{itemize}
  \item Diagnosis and monitoring of gout: acute monoarticular inflammatory arthritis (typically first metatarsophalangeal joint---podagra), with joint aspiration showing negatively birefringent needle-shaped monosodium urate crystals under polarized light microscopy
  \item Monitoring of urate-lowering therapy: target serum urate <6.0 mg/dL for most gout patients; <5.0 mg/dL for patients with tophaceous gout per ACR guidelines
  \item Assessment of tumor lysis syndrome risk in patients initiating chemotherapy for hematologic malignancies (leukemia, lymphoma) and monitoring during treatment. Hyperuricemia >8 mg/dL is a laboratory criterion for TLS per Cairo-Bishop definition
  \item Evaluation of recurrent nephrolithiasis: uric acid stones are radiolucent on plain X-ray, accounting for 5--10\% of all kidney stones, and form in acidic urine (pH <5.5)
  \item Assessment of preeclampsia: hyperuricemia is associated with preeclampsia severity and correlates with poor perinatal outcomes. Uric acid >5.5 mg/dL in third trimester raises concern
  \item Evaluation of suspected Lesch-Nyhan syndrome (X-linked HGPRT deficiency): severe hyperuricemia, self-mutilating behavior, neurological dysfunction in male infants and children
  \item Monitoring of chronic kidney disease: hyperuricemia may be both a marker and a mediator of CKD progression and cardiovascular disease
  \item Screening for inborn errors of purine metabolism in children with neurological symptoms, developmental delay, or unexplained hyperuricemia
\end{itemize}

\section*{Primary vs Secondary Test}
Serum uric acid is a secondary test. It is not part of the standard BMP or CMP and must be specifically ordered. In gout, the diagnosis should ideally be confirmed by joint aspiration and crystal analysis; serum uric acid alone is insufficient because levels may be normal during an acute flare (up to 40\% of patients with acute gout have normal uric acid at the time of the attack due to cytokine-mediated uricosuria and deposition in the affected joint).

\section*{How This Test Is Performed}
A healthcare professional draws blood from a vein into a serum separator tube (gold-top) or lithium heparin tube (green-top). Approximately 3--5 mL of blood is collected using standard venipuncture technique. The sample is centrifuged on an automated chemistry analyzer using the uricase enzymatic method. Results are available within 1--2 hours. No special handling is required beyond standard laboratory practice.

\section*{What Preparation Is Needed}
No fasting is required, though some laboratories recommend an 8-hour fast because recent ingestion of purine-rich foods (red meat, organ meats, shellfish, beer) can transiently elevate uric acid by 1--2 mg/dL. Inform the healthcare provider of all medications: (1) Drugs that elevate uric acid: thiazide and loop diuretics (reduce renal urate excretion, the most common drug cause), low-dose aspirin (<2 g/day inhibits tubular urate secretion), pyrazinamide, ethambutol, cyclosporine, tacrolimus, niacin, levodopa, cytotoxic chemotherapy (tumor lysis). (2) Drugs that lower uric acid: allopurinol, febuxostat, probenecid, benzbromarone, losartan (ARB with mild uricosuric effect), fenofibrate (mild uricosuric effect), high-dose aspirin (>3 g/day), vitamin C (>500 mg/day), corticosteroids, and estrogen (premenopausal women have lower uric acid---postmenopausal levels approach those of men). Alcohol, especially beer (rich in guanosine), increases uric acid production and decreases renal excretion. Fructose-rich beverages increase uric acid production through ATP depletion and purine nucleotide degradation. Recent strenuous exercise increases uric acid from increased ATP turnover. Severe preeclampsia and eclampsia elevate uric acid.

\section*{Contraindications and Precautions}
No absolute contraindications. Important considerations: (1) Uric acid levels are 0.5--1.0 mg/dL higher in males than premenopausal females (estrogen is uricosuric). Postmenopause, female levels approach male levels. (2) Uric acid normally increases with age. (3) Serum uric acid is normal during up to 40\% of acute gout attacks due to acute uricosuria driven by inflammatory cytokines triggering ACTH/cortisol release and due to urate precipitation into the inflamed joint, effectively removing urate from the circulation. Therefore, a normal uric acid does NOT exclude gout and serum urate should be checked 2--4 weeks after the acute flare resolves for a true baseline. (4) Pseudohyperuricemia from ascorbic acid (vitamin C) interference with the uricase-peroxidase method is possible at very high doses. (5) Acetaminophen and salicylates at therapeutic doses may mildly interfere with the peroxidase endpoint method but generally not clinically significant. (6) Bilirubin and hemolysis can interfere with spectrophotometric methods. (7) EDTA or heparinized plasma is acceptable, but EDTA may interfere with some enzymatic methods because EDTA chelates metal cofactors. Serum is preferred.

\section*{Risks}
Minimal risk from venipuncture. Mild pain or bruising resolves in 1--2 days. Rare: hematoma, infection, nerve injury, vasovagal syncope.

\section*{What Happens After the Test}
Results are available within 1--2 hours. Uric acid >6.8 mg/dL exceeds the solubility limit of monosodium urate in plasma at 37 degrees C. However, hyperuricemia alone does not diagnose gout---most people with hyperuricemia never develop gout (only 10--20\% of people with uric acid 7--9 mg/dL develop gout over 10--15 years). Gout diagnosis requires: (1) identification of monosodium urate crystals in synovial fluid by polarized light microscopy (gold standard), OR (2) the ACR/EULAR classification criteria incorporating clinical, laboratory, and imaging findings. Asymptomatic hyperuricemia is generally not treated except in patients with uric acid >13 mg/dL (risk of uric acid nephropathy and tumor lysis syndrome), patients about to receive chemotherapy, or recurrent uric acid stone formers. For gout patients on urate-lowering therapy, target uric acid <6.0 mg/dL (<5.0 for tophaceous gout). In tumor lysis syndrome, uric acid >8 mg/dL with rising trends requires urgent intervention (rasburicase, aggressive IV fluids).

\section*{Understanding Results}

\subsection*{Below Normal}
\begin{itemize}
  \item \textbf{Hypouricemia (<2.0 mg/dL):} Uncommon; most cases due to medications or renal tubular defects.
  \item \textbf{Xanthine oxidase deficiency (hereditary xanthinuria):} Rare autosomal recessive disorder; very low uric acid (<2.0 mg/dL), elevated xanthine and hypoxanthine in urine, xanthine kidney stones. Type I: isolated xanthine oxidase deficiency; Type II: combined xanthine oxidase and aldehyde oxidase deficiency.
  \item \textbf{Uricosuric drugs:} Probenecid, benzbromarone, losartan, fenofibrate, high-dose salicylates, ascorbic acid.
  \item \textbf{SIADH:} Dilutional effect from water retention.
  \item \textbf{Fanconi syndrome:} Generalized proximal tubular dysfunction causing uricosuria, phosphaturia, aminoaciduria, glycosuria. Seen in Wilson disease, cystinosis, heavy metal toxicity, tenofovir, multiple myeloma.
  \item \textbf{Pregnancy:} Increase in GFR increases renal urate clearance, lowering uric acid (2.5--4.5 mg/dL normal in pregnancy).
  \item \textbf{Severe liver disease:} Decreased purine synthesis.
  \item \textbf{Allopurinol or febuxostat therapy:} Intended therapeutic effect; target <6.0 mg/dL.
  \item \textbf{Rasburicase therapy:} Recombinant uricase rapidly degrades uric acid to allantoin; used in tumor lysis syndrome. May cause falsely low uric acid if sample is not immediately placed on ice after collection (ex vivo degradation continues).
\end{itemize}

\subsection*{Above Normal}
\begin{itemize}
  \item \textbf{Hyperuricemia (>6.8 mg/dL in adults):} Defined as concentration above the solubility limit of monosodium urate.
  \item \textbf{Gout (most common clinical consequence):} Recurrent acute inflammatory arthritis from monosodium urate crystal deposition, typically affecting the first MTP joint (podagra), midfoot, ankle, or knee. Chronic tophaceous gout: tophi (aggregates of urate crystals) in joints, bursae, tendons, and soft tissues. Gout is strongly associated with metabolic syndrome, obesity, hypertension, CKD, and cardiovascular disease.
  \item \textbf{Overproduction of uric acid (approximately 10\% of hyperuricemia cases):} Lesch-Nyhan syndrome (HGPRT deficiency, X-linked, severe hyperuricemia, self-mutilation, choreoathetosis, intellectual disability in boys); Kelley-Seegmiller syndrome (partial HGPRT deficiency, milder); PRPP synthetase overactivity (X-linked, early-onset gout and uric acid stones); myeloproliferative disorders (polycythemia vera, essential thrombocythemia, CML) with increased cell turnover; lymphoproliferative disorders; hemolytic anemia; psoriasis (increased epidermal cell turnover); tumor lysis syndrome (massive cell lysis after chemotherapy); high-purine diet (red meat, organ meats, shellfish); fructose-rich beverages; alcohol (especially beer, which contains high guanosine levels); obesity and metabolic syndrome.
  \item \textbf{Underexcretion of uric acid (approximately 90\% of hyperuricemia cases):} Chronic kidney disease (decreased GFR reduces urate filtration). Diuretic therapy (thiazides and loop diuretics are the most common drug cause---volume depletion increases proximal tubular urate reabsorption, and diuretics compete with urate for tubular secretion). Lead nephropathy (saturnine gout---lead toxicity impairs tubular urate secretion; associated with hypertension, CKD, and gout). Lactic acidosis, ketoacidosis, and starvation (ketoacids and lactate compete with urate for tubular secretion). Hypothyroidism. Hyperparathyroidism. Preeclampsia and eclampsia (hyperuricemia >5.5 mg/dL correlates with disease severity). Down syndrome (increased urate production from increased purine turnover). Alcohol (lactic acid production from ethanol metabolism reduces urate excretion). Drugs: cyclosporine, tacrolimus, pyrazinamide, ethambutol, niacin, low-dose aspirin.
  \item \textbf{Tumor lysis syndrome (TLS):} Hyperuricemia >8 mg/dL (or 25\% increase from baseline) plus $\geq$2 of: hyperkalemia >6.0 mEq/L, hyperphosphatemia >6.5 mg/dL, hypocalcemia (corrected Ca <7.0 mg/dL). Prophylaxis with allopurinol 24--48h pre-chemotherapy and aggressive IV hydration (3 L/m\textsuperscript{2}/day). Rasburicase for established TLS with rapidly rising uric acid.
  \item \textbf{Medication-induced hyperuricemia:} Thiazide/loop diuretics, low-dose aspirin, cyclosporine, tacrolimus, pyrazinamide, ethambutol, niacin, levodopa, cytotoxic chemotherapy.
\end{itemize}

\section*{Normal Results}
\begin{itemize}
  \item \textbf{Adult males:} 3.5--7.2 mg/dL (208--428 $\mu$mol/L)
  \item \textbf{Adult females (premenopausal):} 2.6--6.0 mg/dL (155--357 $\mu$mol/L)
  \item \textbf{Adult females (postmenopausal):} 3.5--7.2 mg/dL (similar to males)
  \item \textbf{Children:} 2.0--5.5 mg/dL
  \item \textbf{Pregnancy:} 2.5--4.5 mg/dL
  \item \textbf{Gout treatment target:} <6.0 mg/dL (general); <5.0 mg/dL (tophaceous gout)
  \item \textbf{Solubility threshold for monosodium urate:} 6.8 mg/dL at 37 degrees C and pH 7.4
  \item \textbf{Critical value (TLS risk):} >8 mg/dL or rapid rise
  \item \textbf{Conversion:} 1 mg/dL = 59.5 $\mu$mol/L; 1 $\mu$mol/L = 0.0168 mg/dL
\end{itemize}

\section*{Follow-up Tests}
\begin{itemize}
  \item \textbf{Arthrocentesis with synovial fluid analysis and polarized light microscopy:} Gold standard for gout diagnosis. Negatively birefringent needle-shaped monosodium urate crystals (yellow when parallel to the compensator axis). Rule out septic arthritis with Gram stain and culture.
  \item \textbf{24-hour urine uric acid:} Distinguishes overproducers (>800 mg/24h on normal diet; >600 mg/24h on purine-restricted diet) from underexcretors (<600 mg/24h on normal diet). Guides choice of urate-lowering therapy: underexcretors respond to uricosurics (probenecid); overproducers require xanthine oxidase inhibitors (allopurinol, febuxostat).
  \item \textbf{Renal function (BUN, creatinine, eGFR):} Assess for CKD as cause or consequence of hyperuricemia. Required before starting allopurinol (dose adjustment for eGFR).
  \item \textbf{Complete blood count (CBC):} Assess for myeloproliferative disorders.
  \item \textbf{Serum electrolytes, calcium, phosphate:} Part of TLS monitoring.
  \item \textbf{Urine pH:} Persistently acidic urine (pH <5.5) promotes uric acid stone formation. >24-hour urine with volume, pH, citrate, oxalate, calcium for complete stone risk profile.
  \item \textbf{Radiographs of affected joints:} Chronic gout shows erosions with overhanging edges (Martel sign), joint space preservation until late, soft tissue tophi.
  \item \textbf{Joint ultrasound or dual-energy CT:} Can detect urate crystal deposits non-invasively. Dual-energy CT can differentiate urate from calcium deposition.
  \item \textbf{Serum lead level:} If saturnine gout is suspected (occupational history, hypertension, CKD).
  \item \textbf{HGPRT enzyme assay:} If Lesch-Nyhan syndrome suspected in male children.
\end{itemize}

\section*{References}
\begin{enumerate}
  \item Burtis CA, Ashwood ER, Bruns DE. Tietz Textbook of Clinical Chemistry and Molecular Diagnostics. 7th ed. Elsevier; 2023.
  \item Wallach J. Interpretation of Diagnostic Tests. 10th ed. Wolters Kluwer; 2022.
  \item Tierney LM, Saint S, Drazen JM. CURRENT Medical Diagnosis and Treatment. McGraw-Hill; 2023.
  \item Fischbach FT, Dunning MB. A Manual of Laboratory and Diagnostic Tests. 10th ed. Wolters Kluwer; 2023.
  \item Richette P, Doherty M, Pascual E, et al. 2016 updated EULAR evidence-based recommendations for the management of gout. Ann Rheum Dis. 2017;76(1):29--42.
  \item Neogi T, Jansen TLTA, Dalbeth N, et al. 2015 Gout classification criteria. Ann Rheum Dis. 2015;74(10):1789--1798.
  \item Choi HK, Mount DB, Reginato AM. Pathogenesis of gout. Ann Intern Med. 2005;143(7):499--516.
  \item Cairo MS, Bishop M. Tumour lysis syndrome: new therapeutic strategies and classification. Br J Haematol. 2004;127(1):3--11.
\end{enumerate}

\clearpage
\section*{Regulatory Status and Authorization Date}
Regulatory status applies to the specific assay, instrument, manufacturer, and intended use rather than to the test name alone. No single approval date can be assigned to this test category. For a commercial assay, verify the current product in the relevant regulator's database: the U.S. Food and Drug Administration (FDA) clearance, approval, or authorization; India's Central Drugs Standard Control Organisation (CDSCO) licence or permission; the European Union In Vitro Diagnostic Regulation (EU IVDR) CE certificate issued through the applicable conformity-assessment route; or the United Kingdom Medicines and Healthcare products Regulatory Agency (MHRA) registration and UKCA/CE status. Laboratory-developed tests may not have an individual FDA, CDSCO, EU IVDR, or MHRA approval date.
\clearpage
